\section{Implementation}
\label{sec:impl}

We implement \sys as an end-to-end pipeline in Python. Each of the four agents is realized as an LLM-backed module driven by \texttt{claude-sonnet-5} at temperature 0.7, orchestrated as a LangGraph state machine through a common controller that manages prompting, tool invocation, and the per-agent execute-and-repair loop. Agents communicate through a shared, structured project state---the current source tree, the interface delta emitted by the \emph{LSA}, and the running verification log---serialized as JSON so that any agent can inspect the artifacts produced upstream.

\textbf{Execution and sandboxing.}
Local verification requires actually running untrusted, machine-generated code. \sys stages each candidate clone into an isolated temporary working directory and executes it with the host interpreter under a wall-clock ceiling, capturing compiler diagnostics, exit status, standard output, and any assertion failures. This evidence is returned verbatim to the responsible agent as the repair signal, grounding each retry in concrete runtime feedback rather than speculation. We note the isolation boundary as a limitation rather than a strength: it is filesystem-level, not container-level, so a clone that reached the network or the wider filesystem would not be contained. The corpus is our own and the generated code is inspected, but deployment on untrusted subjects would require a container or VM boundary.

\textbf{Test harness.}
The \emph{VCA} discovers and drives whatever test assets accompany a subject program---unit and integration tests as well as recorded reference input/output traces---through a uniform adapter, so that heterogeneous test formats present a single pass/fail/attribution interface to the pipeline.

\textbf{Equivalence oracle.}
For global verification we use a four-layer source-level oracle, applied in increasing order of cost. The first layer validates that the clone's imports resolve and that the migrated-away distribution is genuinely absent. The second replays recorded golden vectors---inputs and expected outputs captured from the pristine original---at zero numerical tolerance. The third checks program invariants that must hold irrespective of implementation, such as monotonicity of a control output or the reachability of every declared state. The fourth runs the subject's own \texttt{pytest} suite against the clone. A verdict of equivalence requires all four to pass, and the first layer to fail identifies the responsible agent directly.

Evaluation exposed a systematic blind spot in this arrangement, and we added a fifth layer in response. Every subject in the corpus dependency-injects its I/O client so that its tests can run without live hardware, which means the factory function that actually constructs the substituted library is never executed by the suite. The fifth layer calls that factory directly against an unroutable address, and classifies the resulting exception: a connection error is expected and benign, whereas a \texttt{TypeError}, \texttt{AttributeError} or \texttt{ValueError} indicates the migrated call is malformed in a way no absent server could explain. This layer caught defects on clones that had passed all four preceding layers.

\textbf{Toward binary-level confirmation.}
The architecture in Section~\ref{sec:approach} anticipates reusing our prior binary clone detector~\cite{agadakos2025bcd} as the global oracle, disassembling the compiled original and clone, generating emulator execution traces, and prompting an LLM to decide semantic equivalence and justify its verdict. That integration is not exercised in the results reported here; the equivalence-confirmation figures in Section~\ref{sec:results} come from the source-level oracle described above. Coupling the two---so that the detector's localization of divergence is parsed into a targeted repair signal for the responsible agent---remains the most direct extension of this work.

\paragraph{Subject Programs}
Our evaluation corpus consists of twelve Python programs from the industrial measurement and control domain, organized into four categories of increasing transformation demand and ranging in complexity from a single self-contained module to multi-file packages with internal dependencies (Section~\ref{sec:results}, Table~\ref{tab:corpus}). The programs were authored for this study rather than harvested from public repositories, a trade we revisit in Section~\ref{sec:results}: it costs external validity, and it is what makes the per-program golden vectors and invariants the oracle consumes available at all. Every subject ships a \texttt{pytest} suite and recorded input/output traces, which serve both as the \emph{VCA}'s test corpus and as ground truth for the equivalence oracle.
